\subsection{Tree decompositions with logarithmic height}
\label{sec:tw-log-height}

Given a tree decomposition $T$ of logarithmic height, using \autoref{lem:near-opt-1} we are ready to compute a set of partitions $\Pi_B$ of FPT size for each bag $B$, such that a near-optimal solution conforms to $\Pi_B$. In particular, by \autoref{lem:nice-tree-decomp} we may assume that the height of $T$ is $h=O(k\log n)$, which means that the bound on $\zeta_B$ in \autoref{lem:near-opt-1} translates to $O(\frac{k^6}{\eps^2}\log^3\frac{n}{\eps})$.
As in the previous section, we need to apply \autoref{lem:aspect-ratio} in order to bound the aspect ratio of the graph, so that $n$ denotes the number of vertices of the original input graph, while now the graph $G$ has at most $n$ vertices but the ratio between the longest and shortest edge is at most~$2n/\eps$.
We begin by describing how to obtain the near-optimal solution, after which we will identify the partition sets~$\Pi_B$.

\subsubsection{A near-optimal solution}
\label{sec:approx-sol}

\citet{bateni2011approximation} construct a near-optimal solution by modifying the optimum. 
We will use similar techniques to obtain our near-optimal solution, but we construct it by instead modifying the $(1+\eps)$-approx\-imate solution~$F_\eps$ given by \autoref{lem:near-opt-1}.
In particular, we construct a near-optimal $(1+\eps)^2$-approximation~$\widetilde{F}_\eps$ from~$F_\eps$.
The main idea to obtain~$\widetilde{F}_\eps$ is to connect components of~$F_\eps$ if they are very close to one another. 
As before however, doing this naively would incur too much cost for the additional connections.

To make sure that the cost incurred by connecting components of $F_\eps$ is not too large, 
\cite{bateni2011approximation} introduced a partial order on the components based on the structure of a given rooted tree decomposition~$T$.
Let~$C_1,C_2$ be two components of~$F_\eps$ that share a bag $B$ of $T$, i.e.,
$V(C_1)\cap B\neq\emptyset$ and $V(C_2)\cap B\neq\emptyset$.  Since $C_1$ and
$C_2$ are connected subgraphs of the input graph, a basic property of tree
decompositions implies that there are (connected) subtrees $T_1$ and~$T_2$ of
$T$ induced by the respective bags containing vertices of $C_1$ and~$C_2$.
Because these components both contain vertices of~$B$, the node associated with
$B$ is part of both $T_1$ and $T_2$, and therefore the roots of both subtrees
lie on the path from this node to the root of~$T$.  This defines an order on
$C_1$ and $C_2$, and we write $C_1\leq C_2$ if the root of~$T_1$ is farther
from the root of~$T$ than the root of $T_2$ is.  This order is defined for any
two components of~$F_\eps$ that share a bag, and thus we obtain a partial order
on the components of~$F_\eps$, where any components that do not share a bag are
incomparable. 

Using the defined order, \citet{bateni2011approximation} connect components of the optimum solution that are very close to each other.
In order to obtain smaller partition sets, we modify the distance bound used in this procedure compared to \citet{bateni2011approximation}. 
In particular, for any value~$x>0$, let $\pow{x}=2^{\lfloor\log_{2}x\rfloor}$ 
denote the largest power of~$2$ that is at most $x$.
%
Now, starting with $\widetilde{F}_\eps=F_\eps$ we connect components by exhaustively applying the following rule:
\begin{description}
\item[Rule 2:]\customlabel{rule2}{Rule~2} if $C,C'$ are components of $F_\eps$ with $C\leq C'$ and $\dist(C,C')\leq\frac{\eps}{k}\pow{\cost(C)}$ but $C$ and $C'$ lie in different components of $\widetilde{F}_\eps$, then add a shortest path of length $\dist(C,C')$ between $C$ and $C'$ to the solution $\widetilde{F}_\eps$.
\end{description}

A crucial but subtle observation is that for a component~$C$ of $F_\eps$ there can be many components~$C'\leq C$ at distance at most $\frac{\eps}{k}\pow{\cost(C)}$ to $C$, 
which however are not connected to~$C$ in the resulting solution~$\widetilde{F}_\eps$ according to \autoref{rule2}. 
This makes it non-trivial to find small partition sets~$\Pi_B$.
Contrary to this however, an important property of the order on the components is that for any component~$C$ of $F_\eps$, 
there are at most~$k$ other components $C'$ for which~$C\leq C'$, as we will argue for the following lemma to bound the cost of $\widetilde{F}_\eps$. 
In particular, the lemma implies that $\widetilde{F}_\eps$ is a near-optimal $(1+\eps)^2$-approximation, given that~$F_\eps$ is a $(1+\eps)$-approximation.

\begin{lemma}\label{lem:near-opt-2} %
The cost of the solution $\widetilde{F}_\eps$ obtained by \autoref{rule2} from $F_\eps$ is at most $(1+\eps)\cost(F_\eps)$.
\end{lemma}

%
\begin{proof}%
Consider a component $C$ of $F_\eps$ and the highest (closest to the root) node of $T$ for which the bag~$B$ contains a vertex of $C$.
Any component~$C'$ with $C\leq C'$ also intersects $B$, and as this bag
has size at most~$k+1$, there can be at most $k$ such components $C'$.
As a consequence, we can charge the additional cost of connecting a component~$C$ with components $C'$ for which $C\leq C'$
to the cost of~$C$. In particular, if $\mathcal{C}$ denotes the set of all components of $F_\eps$, then
the cost incurred by connecting components according to \autoref{rule2} is at most
\[
\sum_{C\in\mathcal{C}}\sum_{C'\in\mathcal{C}: C\leq C'}\frac{\eps}{k}\pow{\cost(C)}\leq
\sum_{C\in\mathcal{C}}\sum_{C'\in\mathcal{C}: C\leq C'}\frac{\eps}{k}\cost(C)\leq
\sum_{C\in\mathcal{C}}\eps\cost(C)=\eps\cost(F_\eps).
\]
Thus adding the cost of connecting components of $F_\eps$ according to \autoref{rule2}, the cost of the resulting solution is $\cost(\widetilde{F}_\eps)\leq\cost(F_\eps)+\eps\cost(F_\eps)=(1+\eps)\cost(F_\eps)$.
\end{proof}    
%

\subsubsection{Partitioning active terminals}
\label{sec:partitions}

Given the construction of the $(1+\eps)^2$-approximate solution $\widetilde{F}_\eps$ above, the next step is to find a set of partitions $\Pi_B$
of the active terminals $A_B$ for each bag $B$, such that $\widetilde{F}_\eps$ conforms with all sets $\Pi_B$.
In the following, fix a bag $B$ of the given tree decomposition $T$.
The technique used by \citet{bateni2011approximation} is to guess a small net for each active component of bag $B$,\footnote{\citet{bateni2011approximation} refer to these nets as \emph{groups}.}
so that every terminal of $A_B$ close to a net point must be part of the same component in the approximate solution, 
after taking the order on the active components as defined previously into account.
Next we choose a net on the terminals of each active component and bound its size.

\begin{lemma}\label{lem:nets}%
Let $N\subseteq A_B\cap C$ be an $\frac{\eps}{k}\pow{\cost(C)}$-net of the metric induced by 
the active terminals of some active component~$C$.
The size of the net can be bounded by $|N|\leq\lfloor 4k/\eps\rfloor$.\footnote{A slightly worse bound follows from \cite[Lemma~19]{bateni2011approximation}.}
\end{lemma}

%
\begin{proof}%
Let $U$ be a minimum spanning tree of the metric closure of $N$.
It is well known that a minimum spanning tree is a $2(1-\frac{1}{p})$-approximation to an optimum Steiner tree~\cite{KouMB81} on $p$ terminals, 
and thus we have $\cost(C)\geq\frac{1}{2(1-\frac{1}{|N|})}\cdot\cost(U)$, as $C$ in particular is a Steiner tree for $N$. The distance between any pair of 
net points in~$N$ is more than $\frac{\eps}{k}\pow{\cost(C)}\geq \frac{\eps}{2k}\cost(C)$, 
and given that the spanning tree~$U$ has $|N|-1$ edges we get $\cost(U)>\frac{\eps}{2k}\cost(C)(|N|-1)$. 
Putting these two inequalities together we get $\cost(C)>\frac{\eps(|N|-1)}{4k(1-\frac{1}{|N|})}\cost(C)=\frac{\eps|N|}{4k}\cost(C)$, which implies $|N|\leq\lfloor 4k/\eps\rfloor$ as $|N|$ is an integer.
\end{proof}
%

Following the algorithm of \citet{bateni2011approximation}, the next step would be to guess such an 
$\frac{\eps}{k}\pow{\cost(C)}$-net for each of the at most $k+1$ active components $C$ of the bag $B$.
By \autoref{lem:nets}, the total number of net points for these at most $k+1$ nets is at most~$\lfloor 4k/\eps\rfloor(k+1)=O(k^2/\eps)$.
Since however there may be up to~$n$ active terminals, guessing these nets for all active components
can result in $n^{O(k^2/\eps)}$ many possible choices, which leads to an XP time algorithm. 
To circumvent this, we instead consider the partition $\zeta_B$ of the active terminals as given by \autoref{lem:near-opt-1}, 
and guess which of the sets of $\zeta_B$ contains a net point.
We will argue that since the size of $\zeta_B$ is~$O(\frac{k^6}{\eps^2}\log^3\frac{n}{\eps})$ there are only $(\frac{k}{\eps}\log \frac{n}{\eps})^{O(k^2/\eps)}$ possibilities, leading to a faster algorithm.

%
%
%
%
%
%
%
%

More concretely, to compute a set of partitions $\Pi_B$ that $\widetilde{F}_\eps$ conforms to, 
our algorithm considers every sequence $((S_1,\delta_1),(S_2,\delta_2),\ldots,(S_\ell,\delta_\ell),\rho)$
of at most $k+1$ pairs $(S_j,\delta_j)$ and partitions $\rho$ of the index set~$\{1,\ldots,\ell\}$, 
where each $S_j$ is a subset of the parts of~$\zeta_B$ such that $|S_j|\leq \lfloor 4k/\eps\rfloor$, 
and $\delta_j\in\{2^q\mid q\in\mathbb{N}_0\land 0\leq q\leq \log_2(2n^2/\eps)\}$ is an integer power of~$2$ between $1$
and $2n^2/\eps$, where $n$ is the number of vertices of the original input graph in accordance with \autoref{lem:aspect-ratio}.
From every such sequence, the algorithm attempts to construct a partition of the active terminals, and if it
succeeds adds it to the set~$\Pi_B$. As we will show, in this process the algorithm will successfully construct
one partition $\pi$ of $A_B$ that $\widetilde{F}_\eps$ conforms to.

Before describing how a partition of the active terminals arises from such a sequence, 
we bound the number of these sequences, which determines the running time.
By \autoref{lem:near-opt-1}, $|\zeta_B|=O(\frac{k^6}{\eps^2}\log^3\frac{n}{\eps})$ if the tree decomposition~$T$ has logarithmic height, so that there are
at most $\binom{|\zeta_B|}{\lfloor 4k/\eps\rfloor}=(\frac{k}{\eps}\log\frac{n}{\eps})^{O(k/\eps)}$ possible choices for each~$S_j$. Clearly there are $O(\log\frac{n}{\eps})$ choices for each $\delta_j$, and $\ell^\ell=k^{O(k)}$ possible partitions~$\rho$,
given that $\ell\leq k+1$. Since a sequence contains $\ell$ sets $S_j$, the total number of sequences is bounded by~$(\frac{k}{\eps}\log\frac{n}{\eps})^{O(k^2/\eps)}$.

Each sequence may give rise to a partition $\pi\in\Pi_B$ of the active terminals as follows.
First, let $\pi=\{Y_1,\ldots,Y_{|\rho|}\}$, i.e., $\pi$ has the same number of sets as the partition $\rho$.
Let $U_j=\bigcup_{U\in S_j} U$ denote the set of active terminals in~$S_j$, 
and let $\rho(j)$ be the part of $\rho$ containing $j$.
We distinguish between active terminals $t\in A_B$ that lie in some set $U_j$ and those that do not:
\begin{itemize}
    \item if $t\in U_j$ for some $j\in[\ell]$ then $t\in Y_{\rho(j)}$ (i.e., $U_j\subseteq Y_{\rho(j)}$), and
    \item otherwise, if $p_t\in\{1,\ldots,\ell\}$ denotes the smallest index for which $\dist(t,U_{p_t})\leq\frac{\eps}{k}\delta_{p_t}$, then $t\in Y_{\rho(p_t)}$.
\end{itemize}
If this $\pi$ is a partition of $A_B$ we add $\pi$ to $\Pi_B$, and otherwise we dismiss the current sequence.
Clearly $\pi$ can be constructed in polynomial time, given a sequence.

\begin{lemma}
The $(1+\eps)^2$-approximate solution~$\widetilde{F}_\eps$ conforms to the set $\Pi_B$ of partitions constructed above.
\end{lemma}
\begin{proof}
Consider the $(1+\eps)$-approximate solution $F_\eps$ of \autoref{lem:near-opt-1} from which $\widetilde{F}_\eps$ is constructed 
according to \autoref{rule2}, and the partition $\zeta_B$ of $A_B$ as given by \autoref{lem:near-opt-1}. Let the active components of 
$F_\eps$ be $C_1,\ldots,C_\ell$ indexed according to their order, i.e., $C_j\leq C_{j'}$ if and only if~$j\leq  j'$.
For each active component~$C_j$ we fix an $\frac{\eps}{k}\pow{\cost(C_j)}$-net $N_j$ of size at most~$\lfloor 4k/\eps\rfloor$
according to \autoref{lem:nets}. Now, consider the sequence $((S_1,\delta_1),(S_2,\delta_2),\ldots,(S_\ell,\delta_\ell),\rho)$, where
\begin{itemize}
    \item $S_j$ contains exactly those sets of $\zeta_B$ that contain at least one net point of $N_j$,
    \item $\delta_j=\pow{\cost(C_j)}$, and
    \item $\rho$ is the partition of the index set corresponding to the components of $\widetilde{F}_\eps$, i.e., $\rho(j)=\rho(j')$ if and only if $C_j$ and $C_{j'}$ lie in the same component in $\widetilde{F}_\eps$.
\end{itemize}

Recall that after applying \autoref{lem:aspect-ratio} to the input, the ratio between the shortest and longest edge 
is at most $2n/\eps$, where $n$ is the number of vertices of the original input graph.
Since we assume that the length of the shortest edge is $1$, the cost of any component lies between $1$ and~$2n^2/\eps$, 
given that a component is a tree with less than $n$ edges. Therefore $\pow{\cost(C_j)}\in\{2^q\mid q\in\mathbb{N}_0\land 0\leq q\leq \log_2(2n^2/\eps)\}$, which means that the algorithm will consider the above sequence in some iteration.

We now turn to $\pi=\{Y_1,\ldots,Y_{|\rho|}\}$ constructed for this sequence, and show that it is a partition of~$A_B$ and that $\widetilde{F}_\eps$ conforms to it.
For this, note that no set of $\zeta_B$ contains net points of several active components of $F_\eps$,
since by \autoref{lem:near-opt-1} all active terminals in the same set of $\zeta_B$ also belong to the same component of $F_\eps$.
Thus the sets $S_j$ as defined above (and also the corresponding sets $U_j$) are pairwise disjoint. 
This means that, due to the definition of $\rho$, any two terminals $t\in U_j$ and $t'\in U_{j'}$ end up in the same set
of $\pi$ if and only if $t$ and $t'$ belong to the same component of $\widetilde{F}_\eps$ (as $U_j\subseteq Y_{\rho(j)}$).

Now consider a terminal $t\in A_B$, which does not lie in any $U_j$, and let $q$ be the index of the active 
component $C_q$ of $F_\eps$ containing~$t$. As $\delta_q=\pow{\cost(C_q)}$, $N_q$ is an $\frac{\eps}{k}\delta_q$-net of 
$C_q\cap A_B$. Also, we chose $S_q$ so that $N_q\subseteq U_q$. Hence we get 
$\dist(t,U_q)\leq\dist(t,N_q)\leq\frac{\eps}{k}\delta_q$, and the definition of $p_t$ implies $p_t\leq q$.
Now~$C_{p_t}$ is either equal to~$C_q$, or~$C_q$ is connected to the component~$C_{p_t}$ in the approximate solution~$\widetilde{F}_\eps$
according to \autoref{rule2}: on one hand we have $C_{p_t}\leq C_q$ due to the order of the indices, and at the same time by 
\autoref{lem:near-opt-1} we have $U_{p_t}\subseteq V(C_{p_t})\cap A_B$, which implies
\[
\dist(C_q,C_{p_t})\leq\dist(t,C_{p_t})\leq\dist(t,U_{p_t})\leq\frac{\eps}{k}\delta_{p_t}=\frac{\eps}{k}\pow{\cost(C_{p_t})}.
\]
Hence we can conclude that~$t$ lies in the same component as $C_{p_t}$ in $\widetilde{F}_\eps$. 

In conclusion, adding $U_j$ to $Y_{\rho(j)}$ and~$t$ to $Y_{\rho(p_t)}$ for each terminal $t$ not lying in any $U_j$, 
partitions the terminals according to the components of $\widetilde{F}_\eps$.
Hence $\pi$ is a partition of the active terminals~$A_B$ that is added to $\Pi_B$, and~$\widetilde{F}_\eps$ conforms to it.
\end{proof}

Using all of the above, we can finally prove our main theorem, stating that there is an EPAS for \SF parameterized by the treewidth.

\begin{proof}[Proof of \autoref{thm:tw-EPAS}]

The first steps of our algorithm are to preprocess the given tree decomposition
using \autoref{lem:nice-tree-decomp} so that it is nice and its height is
$O(k\log n)$, and the input graph using \autoref{lem:aspect-ratio} so that the
aspect ratio is bounded (which means that $n$ denotes the number of vertices in
the original input graph). We then compute the partition sets $\Pi_B$ for all
bags $B$ using the above procedure, resulting in partition sets of size
$(\frac{k}{\eps}\log
\frac{n}{\eps})^{O(k^2/\eps)}=2^{O(\frac{k^2}{\eps}\log\frac{k}{\eps})}\cdot
n^{o(1)}$. Here, we are using a well-known Win/Win argument: if
$k^2/\eps<\sqrt{\log n}$, then $(\log n)^{k^2/\eps} = n^{o(1)}$; otherwise,
$\log n\le k^4/\eps^2$, therefore  $(\frac{k}{\eps}\log
\frac{n}{\eps})^{O(k^2/\eps)}=(\frac{k}{\eps})^{O(\frac{k^2}{\eps})}$.


Since each partition of a set $\Pi_B$ can be computed in polynomial time, and
the number of bags of the nice tree decomposition is $O(kn)$, this takes
$2^{O(\frac{k^2}{\eps}\log\frac{k}{\eps})}\cdot n^{O(1)}$ time. Next we apply
\autoref{thm:tw-dp} to compute a solution that is at least as good as
$\widetilde{F}_\eps$ conforming to all $\Pi_B$, in
$2^{O(\frac{k^2}{\eps}\log\frac{k}{\eps})}\cdot n^{O(1)}$ time. Hence we obtain
a $(1+\eps)^2$-approximation~$F$. According to \autoref{lem:aspect-ratio}, $F$
can be converted into a $((1+\eps)^2+\eps)$-approximation to the original input
graph. Since for any $\eps'>0$ we may choose $\eps=\Theta(\eps')$ so that
$((1+\eps)^2+\eps)\leq 1+\eps'$, we obtain an EPAS as claimed.  \end{proof}
